\SourcePage{355}{360}
\section{Scattering of an electron in an external field}

Elastic scattering of an electron in a constant external field represents the simplest process, existing already in the first approximation of perturbation theory (first Born approximation). To it there corresponds a diagram with one vertex
% [Feynman diagram (80.1): Single vertex with external field (wavy line with cross, momentum $k$) and two electron lines (momenta $p$ and $p'$)]
\begin{equation}
\tag{80.1}
\end{equation}
where $p$ and $p'$ are the initial and final 4-momenta of the electron, and $q = p' - p$. Since the energy of the electron is conserved in scattering in a constant field ($\varepsilon = \varepsilon'$), we have $q = (0, \mathbf{q})$\,\footnote{In the case of an external field such a diagram is of course not forbidden by the law of conservation of 4-momentum (as it was in diagram~(73.19) with a real photon): the square $q^2$, in contrast to the square of the 4-momentum of a real photon, need not be equal to zero; from the Fourier integral representing the external field, the component with the required~$q$ is automatically selected.}.

The corresponding scattering amplitude
\begin{equation}
M_{fi} = -e\bar{u}(p')[\gamma A^{(e)}(\mathbf{q})]u(p),
\tag{80.2}
\end{equation}
where $A^{(e)}(\mathbf{q})$ is the component of the spatial Fourier expansion of the external field. The scattering cross section, according to~(64.26),
\begin{equation}
d\sigma = \frac{1}{16\pi^2}|M_{fi}|^2\,do'.
\tag{80.3}
\end{equation}

For an electrostatic field $A^{(e)} = (A_0^{(e)}, 0)$, so that
\begin{equation}
M_{fi} = -e\bar{u}(p')\gamma^0 u(p)A_0^{(e)}(\mathbf{q}) = -eu^*(p')u(p)A_0^{(e)}(\mathbf{q}).
\tag{80.4}
\end{equation}
In the non-relativistic case the bispinor amplitudes of the plane waves $u(p)$ reduce to the non-relativistic (two-component) amplitudes. For scattering without change of polarisation this is a quantity independent of $p$, and by virtue of the normalisation condition we have adopted $u^*u = 2m$. Taking this into account, we obtain
\[
d\sigma = \left|-\frac{m}{2\pi}U(\mathbf{q})\right|^2 do',
\]

\SourcePage{356}{361}
where $U(\mathbf{q}) = eA_0^{(e)}(\mathbf{q})$ is the component of the Fourier expansion of the potential energy of the electron in the field; this expression coincides with the well-known Born formula (see~III,~(126.7)).

In the general relativistic case the scattering cross section of unpolarised electrons is obtained by averaging the square of the modulus $|M_{fi}|^2$ over the initial and summing over the final polarisations, i.e.\ by forming the quantity
\[
\frac{1}{2}\sum_{\text{polar}}|M_{fi}|^2,
\]
where the summation is carried out over the spin directions of the initial and final electrons; the factor $1/2$ converts one of these summations into an averaging. By the rules set out in~\S\,65 we obtain
\[
\frac{1}{2}\sum_{\text{polar}}|M_{fi}|^2 = 2\operatorname{Sp}\rho'(\gamma A_0^{(e)}) \rho (\gamma A_0^{(e)*}) =
\]
\[
= \frac{1}{2}|A_0^{(e)}(\mathbf{q})|^2\operatorname{Sp}(m + \gamma p')\gamma^0(m + \gamma p)\gamma^0.
\]
For computing the trace we note that $\gamma^0(\gamma p)\gamma^0 = \gamma\widetilde{p}$, where $\widetilde{p} = (\varepsilon, -\mathbf{p})$, and therefore
\[
\frac{1}{4}\operatorname{Sp}(m + \gamma p')\gamma^0(m + \gamma p)\gamma^0 = \frac{1}{4}\operatorname{Sp}(m + \gamma p')(m + \gamma\widetilde{p}) =
\]
\[
= m^2 + p'\widetilde{p} = \varepsilon^2 + m^2 + \mathbf{p}\mathbf{p}' = 2\varepsilon^2 - \mathbf{q}^2/2.
\]

Hence the cross section
\begin{equation}
d\sigma = \frac{e^2|A_0^{(e)}(\mathbf{q})|^2}{4\pi^2}\varepsilon^2\left(1 - \frac{\mathbf{q}^2}{4\varepsilon^2}\right)do'.
\tag{80.5}
\end{equation}

For the field created by a static distribution of charges with density $\rho(\mathbf{r})$, we have
\begin{equation}
A_0^{(e)}(\mathbf{q}) = \frac{4\pi\rho(\mathbf{q})}{\mathbf{q}^2},
\tag{80.6}
\end{equation}
where $\rho(\mathbf{q})$ is the Fourier transform of the distribution $\rho(\mathbf{r})$ (form factor). In particular, for the Coulomb field of a point charge $Ze$ we have: $\rho(\mathbf{q}) = Ze$. Then the scattering cross section
\begin{equation}
d\sigma = do'\,\frac{4(Ze^2)^2\varepsilon^2}{\mathbf{q}^4}\left(1 - \frac{\mathbf{q}^2}{4\varepsilon^2}\right)
\tag{80.7}
\end{equation}

\SourcePage{357}{362}
(\textit{N.\,F.\,Mott}, 1929). The square
\[
\mathbf{q}^2 = 4\mathbf{p}^2\sin^2(\theta/2),
\]
where $\theta$ is the scattering angle. Therefore the expression in front of the bracket by its angular dependence may be called the Rutherford cross section:
\begin{equation}
d\sigma_\text{Ru} = do\,\frac{4(Ze^2)^2\varepsilon^2}{\mathbf{q}^4} = do\,\frac{4(Ze^2)^2\varepsilon^2}{4\mathbf{p}^4}\sin^{-4}\frac{\theta}{2}
\tag{80.8}
\end{equation}
(in the non-relativistic limit the coefficient $\varepsilon^2/\mathbf{p}^4 \to 1/(m^2 v^4)$).
Thus\,\footnote{The fact that this formula expresses the difference of $d\sigma$ from $d\sigma_\text{Ru}$ is specific to particles with spin~$1/2$. For the scattering of particles with spin~0 (if their motion in the electromagnetic field were described by the wave equation) we would obtain $d\sigma = d\sigma_\text{Ru}$. At first glance it appears strange that this purely quantum effect does not contain~$\hbar$. One must, however, keep in mind that the condition of applicability of the Born approximation ($e^2/(\hbar v) \ll 1$) is the opposite of the condition of quasiclassicality for the motion in a Coulomb field and therefore the passage to the classical case in formula~(80.9) is impossible.},
\begin{equation}
d\sigma = d\sigma_\text{Ru}\left(1 - v^2\sin^2\frac{\theta}{2}\right).
\tag{80.9}
\end{equation}

Let us note that in the ultrarelativistic case the angular distribution differs from the non-relativistic one by a strong suppression of backward scattering (when $\theta \to \pi$: $d\sigma/d\sigma_\text{Ru} \to m^2/\varepsilon^2$).

In the ultrarelativistic case for scattering at small angles (80.7) gives
\begin{equation}
d\sigma = \frac{4(Ze^2)^2}{\varepsilon^4\theta^4}\,do'.
\tag{80.10}
\end{equation}

Although we have obtained this formula in the Born approximation (i.e.\ assuming $Ze^2 \ll 1$), it nevertheless remains valid (for angles $\theta \lesssim m/\varepsilon$) also when $Ze^2 \sim 1$. One can verify this with the help of the ultrarelativistic exact (in $Ze^2$) wave function $\psi_{\mathbf{p}}^{(+)}$~(39.10). This solution, valid in the region~(39.2), remains, of course, valid also in the asymptotic region for arbitrarily large~$r$. Here
\[
F \propto 1 + \text{const}\cdot e^{i(pr-\mathbf{pr})},\qquad \frac{\alpha\nabla F}{\varepsilon} \sim 1 - \cos\theta \sim \theta^2 \ll 1,
\]
so that the correction term remains, as one would expect, small. The wave function is of the type $e^{i\mathbf{pr}}F$, coinciding in form with the non-relativistic function (with an obvious change of parameters), has the same asymptotic form, and therefore also for the cross section one obtains the Rutherford expression.

For computing the cross section of scattering of arbitrarily polarised electrons one could use the general rules of density matrices (see~\S\,29.13). In the present case, however,

\SourcePage{358}{363}
the result can be obtained in a less cumbersome way, by representing the bispinor amplitudes $u(p')$ and $u(p)$ in the form~(23.9); multiplying them out, we obtain
\[
u^*(p')u(p) = w'^*\{\varepsilon + m + (\varepsilon - m)(\mathbf{n}'\boldsymbol{\sigma})(\mathbf{n}\boldsymbol{\sigma})\}w,
\]
or, using formula~(33.5),
\begin{equation}
u^*(p')u(p) = w'^*\hat{f}w,
\tag{80.11}
\end{equation}
where\,\footnote{The definitions of $\hat{f}$ here and in~\S\,37, \S\,38 differ by a common factor.}
\begin{equation}
\hat{f} = A + B\boldsymbol{\nu}\boldsymbol{\sigma},
\tag{80.12}
\end{equation}
\[
A = (\varepsilon + m) + (\varepsilon - m)\cos\theta,\quad B = -i(\varepsilon - m)\sin\theta,\quad \boldsymbol{\nu} = \frac{[\mathbf{n}\mathbf{n}']}{{\sin\theta}}.
\]

The two-component quantity (3-spinor) $w$ represents the non-relativistic spinor wave function of the electron. The transition to partially polarised states is accomplished by replacing products $w_\alpha w^*_\beta$ ($\alpha$, $\beta$ are spinor indices) by the non-relativistic two-row density matrix $\rho_{\alpha\beta}$.
Thus, one must replace
\[
|M_{fi}|^2 \to e^2|A_0^{(e)}(\mathbf{q})|^2\operatorname{Sp}\rho'(A - B\boldsymbol{\nu}\boldsymbol{\sigma})\rho'(A + B\boldsymbol{\nu}\boldsymbol{\sigma}),
\]
where
\[
\rho = \tfrac{1}{2}(1 + \boldsymbol{\sigma}\boldsymbol{\zeta}),\qquad \rho' = \tfrac{1}{2}(1 + \boldsymbol{\sigma}\boldsymbol{\zeta}'),
\]
and $\boldsymbol{\zeta}$ and $\boldsymbol{\zeta}'$ are the vectors of the initial and final polarisations, the latter selected by a detector. Computing the trace leads to the result
\begin{equation}
d\sigma = d\sigma_0\left\{1 + \frac{(A^2 - |B|^2)\boldsymbol{\zeta}\boldsymbol{\zeta}' + 2|B|^2(\boldsymbol{\nu}\boldsymbol{\zeta})(\boldsymbol{\nu}\boldsymbol{\zeta}') + 2A|B|[\boldsymbol{\nu}\boldsymbol{\zeta}\boldsymbol{\zeta}']}{A^2 + |B|^2}\right\},
\tag{80.13}
\end{equation}
where $d\sigma_0$ is the cross section of scattering of unpolarised electrons.

Representing the curly brackets in~(80.13) in the form $\{1 + \boldsymbol{\zeta}^{(f)}\boldsymbol{\zeta}'\}$, we find the polarisation of the final electron as such (in contrast to the detectable final polarisation $\boldsymbol{\zeta}'$---see~\S\,65)\,\footnote{Formula~(80.14) corresponds to the formula found in Problem~1, (see~III, \S\,140) and is obtained from it with real $A$ and imaginary~$B$.}:
\begin{equation}
\boldsymbol{\zeta}^{(f)} = \frac{(A^2 - |B|^2)\boldsymbol{\zeta} + 2|B|^2(\boldsymbol{\nu}\boldsymbol{\zeta})\boldsymbol{\nu} + 2A|B|[\boldsymbol{\nu}\boldsymbol{\zeta}]}{A^2 + |B|^2}.
\tag{80.14}
\end{equation}

We see that the scattered electrons are polarised only if the incident electrons are polarised. This is a general property of the first Born approximation (cf.~III, \S\,140).

\SourcePage{359}{364}
In the non-relativistic case ($\varepsilon \to m$) from~(80.14) we obtain $\boldsymbol{\zeta}^{(f)} = \boldsymbol{\zeta}$, i.e.\ the electron preserves its polarisation upon scattering (a natural consequence of neglecting spin-orbit interaction).

In the opposite, ultrarelativistic, case we have
\[
A = \varepsilon(1 + \cos\theta),\qquad B = -i\varepsilon\sin\theta
\]
(in accordance with the general formula~(38.2)).

If furthermore the incident electron has a definite helicity ($\boldsymbol{\zeta} = 2\lambda\mathbf{n}$, $\lambda = \pm 1/2$), then from~(80.14) it follows after simple reduction that
\[
\boldsymbol{\zeta}^{(f)} = 2\lambda\mathbf{n}'.
\]
In other words, after scattering the electron remains helical, preserving the former value ($\lambda$) of the helicity.

This property, as was already noted in~\S\,38, is connected with the fact that on neglecting mass the Dirac equation in the spinor representation separates into two independent equations for the functions~$\xi$ and $\eta$. This result has also a more general significance, since the current
\[
\hat{j} = (\xi^*\xi + \eta^*\eta,\, \xi^*\boldsymbol{\sigma}\xi - \eta^*\boldsymbol{\sigma}\eta),
\]
and with it the operator of electromagnetic perturbation $\hat{V} = e\hat{j}\hat{A}$, does not contain mixed terms in $\xi$ and $\eta$, and therefore has no matrix elements for transitions between $\xi$- and $\eta$-states. From this it follows that if an ultrarelativistic electron has a definite helicity (i.e.\ differs from zero either $\xi$ or $\eta$), then in processes of interaction this helicity will be conserved in the approximation corresponding to the complete neglect of the electron mass.
